\documentclass[a4paper,12pt]{article}
\usepackage[utf8]{inputenc}
\usepackage[T1]{fontenc}
\usepackage[italian]{babel}
\usepackage{pgfplots}
\usepackage{amsmath}
\usepackage{float}
\usepackage{geometry}
\geometry{margin=2cm}

\pgfplotsset{compat=1.18}

\title{Report 4.4 -- Combinazioni di k e TIMES}

\begin{document}

\maketitle

\newpage

\section{Grafici dei risultati per $k$ e TIMES}

I grafici mostrano rispettivamente l'andamento di:
\begin{itemize}
    \item Avg Profit GAP
    \item Avg Weight GAP
    \item Avg fQ
    \item fQ Min Found
    \item Profit Found
    \item Weight Found
    \item Profit GAP
    \item Weight GAP
    \item n of fQ Min Found
\end{itemize}
al variare di TIMES per diversi valori di $k \in \{1,10,100\}$, confrontando
i risultati usando QALS e non utilizzando QALS.

I risultati riportati sono stati ottenuti utilizzando:

\begin{itemize}
    \item $\lambda = \frac{\sum_{i=1}^{n} p_i}{650 \cdot C}$
    \item l'istanza \texttt{ksp\_2.txt}
\end{itemize}

\begin{figure}[H]
\centering
\begin{tikzpicture}
\begin{axis}[
    width=0.95\textwidth,
    height=0.48\textwidth,
    xmode=log,
    log basis x=10,
    xlabel={TIMES},
    ylabel={Avg Profit GAP},
    title={Avg Profit GAP},
    grid=both,
    legend style={at={(0.5,-0.25)}, anchor=north, legend columns=3},
    xtick={1,10,100,1000},
    xticklabels={1,10,100,1000},
    tick label style={font=\small},
    label style={font=\small},
    title style={font=\small},
]
\addplot[blue, thick, mark=*] coordinates {(1, -4539) (10, -5166.3) (100, -5344.6) (1000, -5232.7)};
\addplot[red, thick, mark=square*] coordinates {(1, -3658) (10, -4702.3) (100, -4915.2) (1000, -5311.5)};
\addplot[green!60!black, thick, mark=triangle*] coordinates {(1, -3082) (10, -5045.3) (100, -5186.5) (1000, -5352)};
\addplot[blue, thick, dashed, mark=o] coordinates {(1, 1528) (10, 676.5) (100, 925.9) (1000, 802.3)};
\addplot[red, thick, dashed, mark=square] coordinates {(1, 0) (10, 24.2) (100, 76.1) (1000, 80.2)};
\addplot[green!60!black, thick, dashed, mark=triangle] coordinates {(1, 0) (10, 0) (100, 0) (1000, 0)};
\legend{
    QALS ($k=1$),
    QALS ($k=10$),
    QALS ($k=100$),
    NO QALS ($k=1$),
    NO QALS ($k=10$),
    NO QALS ($k=100$)
}
\end{axis}
\end{tikzpicture}
\caption{Andamento di \textit{Avg Profit GAP} al variare di TIMES per diversi valori di $k$.}
\label{fig:avg-profit-gap}
\end{figure}

\begin{figure}[H]
\centering
\begin{tikzpicture}
\begin{axis}[
    width=0.95\textwidth,
    height=0.48\textwidth,
    xmode=log,
    log basis x=10,
    xlabel={TIMES},
    ylabel={Avg Weight GAP},
    title={Avg Weight GAP},
    grid=both,
    legend style={at={(0.5,-0.25)}, anchor=north, legend columns=3},
    xtick={1,10,100,1000},
    xticklabels={1,10,100,1000},
    tick label style={font=\small},
    label style={font=\small},
    title style={font=\small},
]
\addplot[blue, thick, mark=*] coordinates {(1, -8233) (10, -8038.2) (100, -8482.5) (1000, -8399.5)};
\addplot[red, thick, mark=square*] coordinates {(1, -7658) (10, -8452.3) (100, -8298.4) (1000, -8393.5)};
\addplot[green!60!black, thick, mark=triangle*] coordinates {(1, -9475) (10, -8105.9) (100, -8256.5) (1000, -8452.9)};
\addplot[blue, thick, dashed, mark=o] coordinates {(1, -54) (10, -94.8) (100, -65.1) (1000, -73.3)};
\addplot[red, thick, dashed, mark=square] coordinates {(1, 0) (10, -16.3) (100, -28.6) (1000, -32.4)};
\addplot[green!60!black, thick, dashed, mark=triangle] coordinates {(1, 0) (10, 0) (100, 0) (1000, 0)};
\legend{
    QALS ($k=1$),
    QALS ($k=10$),
    QALS ($k=100$),
    NO QALS ($k=1$),
    NO QALS ($k=10$),
    NO QALS ($k=100$)
}
\end{axis}
\end{tikzpicture}
\caption{Andamento di \textit{Avg Weight GAP} al variare di TIMES per diversi valori di $k$.}
\label{fig:avg-weight-gap}
\end{figure}

\begin{figure}[H]
\centering
\begin{tikzpicture}
\begin{axis}[
    width=0.95\textwidth,
    height=0.48\textwidth,
    xmode=log,
    log basis x=10,
    xlabel={TIMES},
    ylabel={Avg fQ},
    title={Avg fQ},
    grid=both,
    legend style={at={(0.5,-0.25)}, anchor=north, legend columns=3},
    xtick={1,10,100,1000},
    xticklabels={1,10,100,1000},
    tick label style={font=\small},
    label style={font=\small},
    title style={font=\small},
]
\addplot[blue, thick, mark=*] coordinates {(1, 3104583.38) (10, 3025645.33) (100, 3384450.39) (1000, 3312842.51)};
\addplot[red, thick, mark=square*] coordinates {(1, 2677807.06) (10, 3305083.47) (100, 3225368.19) (1000, 3305006.71)};
\addplot[green!60!black, thick, mark=triangle*] coordinates {(1, 4135708.56) (10, 3093711.48) (100, 3218387.32) (1000, 3355755.13)};
\addplot[blue, thick, dashed, mark=o] coordinates {(1, -53200.03) (10, -53825.79) (100, -53687.67) (1000, -53778.81)};
\addplot[red, thick, dashed, mark=square] coordinates {(1, -54748.31) (10, -54689.85) (100, -54608.46) (1000, -54594.86)};
\addplot[green!60!black, thick, dashed, mark=triangle] coordinates {(1, -54748.31) (10, -54748.31) (100, -54748.31) (1000, -54748.31)};
\legend{
    QALS ($k=1$),
    QALS ($k=10$),
    QALS ($k=100$),
    NO QALS ($k=1$),
    NO QALS ($k=10$),
    NO QALS ($k=100$)
}
\end{axis}
\end{tikzpicture}
\caption{Andamento di \textit{Avg fQ} al variare di TIMES per diversi valori di $k$.}
\label{fig:avg-fq}
\end{figure}

\begin{figure}[H]
\centering
\begin{tikzpicture}
\begin{axis}[
    width=0.95\textwidth,
    height=0.48\textwidth,
    xmode=log,
    log basis x=10,
    xlabel={TIMES},
    ylabel={fQ Min Found},
    title={fQ Min Found},
    grid=both,
    legend style={at={(0.5,-0.25)}, anchor=north, legend columns=3},
    xtick={1,10,100,1000},
    xticklabels={1,10,100,1000},
    tick label style={font=\small},
    label style={font=\small},
    title style={font=\small},
]
\addplot[blue, thick, mark=*] coordinates {(1, 3104583.38) (10, 1947492.7) (100, 1428869.96) (1000, 857615.45)};
\addplot[red, thick, mark=square*] coordinates {(1, 2677807.06) (10, 2212399.92) (100, 1692743.31) (1000, 1289900.58)};
\addplot[green!60!black, thick, mark=triangle*] coordinates {(1, 4135708.56) (10, 1756229.12) (100, 1283760.22) (1000, 782294.93)};
\addplot[blue, thick, dashed, mark=o] coordinates {(1, -53200.03) (10, -54431.03) (100, -54748.31) (1000, -54748.31)};
\addplot[red, thick, dashed, mark=square] coordinates {(1, -54748.31) (10, -54748.31) (100, -54748.31) (1000, -54748.31)};
\addplot[green!60!black, thick, dashed, mark=triangle] coordinates {(1, -54748.31) (10, -54748.31) (100, -54748.31) (1000, -54748.31)};
\legend{
    QALS ($k=1$),
    QALS ($k=10$),
    QALS ($k=100$),
    NO QALS ($k=1$),
    NO QALS ($k=10$),
    NO QALS ($k=100$)
}
\end{axis}
\end{tikzpicture}
\caption{Andamento di \textit{fQ Min Found} al variare di TIMES per diversi valori di $k$.}
\label{fig:fq-min-found}
\end{figure}

\begin{figure}[H]
\centering
\begin{tikzpicture}
\begin{axis}[
    width=0.95\textwidth,
    height=0.48\textwidth,
    xmode=log,
    log basis x=10,
    xlabel={TIMES},
    ylabel={Profit Found},
    title={Profit Found: Profitto soluzione con fq minima trovata},
    grid=both,
    legend style={at={(0.5,-0.25)}, anchor=north, legend columns=3},
    xtick={1,10,100,1000},
    xticklabels={1,10,100,1000},
    tick label style={font=\small},
    label style={font=\small},
    title style={font=\small},
]
\addplot[blue, thick, mark=*] coordinates {(1, 12279) (10, 11298) (100, 11396) (1000, 8644)};
\addplot[red, thick, mark=square*] coordinates {(1, 11398) (10, 13071) (100, 11416) (1000, 11108)};
\addplot[green!60!black, thick, mark=triangle*] coordinates {(1, 10822) (10, 9239) (100, 11720) (1000, 11345)};
\addplot[blue, thick, dashed, mark=o] coordinates {(1, 6212) (10, 7443) (100, 7740) (1000, 7740)};
\addplot[red, thick, dashed, mark=square] coordinates {(1, 7740) (10, 7740) (100, 7740) (1000, 7740)};
\addplot[green!60!black, thick, dashed, mark=triangle] coordinates {(1, 7740) (10, 7740) (100, 7740) (1000, 7740)};
\legend{
    QALS ($k=1$),
    QALS ($k=10$),
    QALS ($k=100$),
    NO QALS ($k=1$),
    NO QALS ($k=10$),
    NO QALS ($k=100$)
}
\end{axis}
\end{tikzpicture}
\caption{Andamento di \textit{Profit Found} al variare di TIMES per diversi valori di $k$.}
\label{fig:profit-found}
\end{figure}

\begin{figure}[H]
\centering
\begin{tikzpicture}
\begin{axis}[
    width=0.95\textwidth,
    height=0.48\textwidth,
    xmode=log,
    log basis x=10,
    xlabel={TIMES},
    ylabel={Weight Found},
    title={Weight Found: Peso soluzione con fq minima trovata},
    grid=both,
    legend style={at={(0.5,-0.25)}, anchor=north, legend columns=3},
    xtick={1,10,100,1000},
    xticklabels={1,10,100,1000},
    tick label style={font=\small},
    label style={font=\small},
    title style={font=\small},
]
\addplot[blue, thick, mark=*] coordinates {(1, 9211) (10, 7538) (100, 6630) (1000, 5412)};
\addplot[red, thick, mark=square*] coordinates {(1, 8636) (10, 7959) (100, 7109) (1000, 6360)};
\addplot[green!60!black, thick, mark=triangle*] coordinates {(1, 10453) (10, 7215) (100, 6349) (1000, 5233)};
\addplot[blue, thick, dashed, mark=o] coordinates {(1, 1032) (10, 1032) (100, 978) (1000, 978)};
\addplot[red, thick, dashed, mark=square] coordinates {(1, 978) (10, 978) (100, 978) (1000, 978)};
\addplot[green!60!black, thick, dashed, mark=triangle] coordinates {(1, 978) (10, 978) (100, 978) (1000, 978)};
\legend{
    QALS ($k=1$),
    QALS ($k=10$),
    QALS ($k=100$),
    NO QALS ($k=1$),
    NO QALS ($k=10$),
    NO QALS ($k=100$)
}
\end{axis}
\end{tikzpicture}
\caption{Andamento di \textit{Weight Found} al variare di TIMES per diversi valori di $k$.}
\label{fig:weight-found}
\end{figure}

\begin{figure}[H]
\centering
\begin{tikzpicture}
\begin{axis}[
    width=0.95\textwidth,
    height=0.48\textwidth,
    xmode=log,
    log basis x=10,
    xlabel={TIMES},
    ylabel={Profit GAP},
    title={Profit GAP: Differenza tra la soluzione con fq minima trovata e la soluzione di DP},
    grid=both,
    legend style={at={(0.5,-0.25)}, anchor=north, legend columns=3},
    xtick={1,10,100,1000},
    xticklabels={1,10,100,1000},
    tick label style={font=\small},
    label style={font=\small},
    title style={font=\small},
]
\addplot[blue, thick, mark=*] coordinates {(1, -4539) (10, -3558) (100, -3656) (1000, -904)};
\addplot[red, thick, mark=square*] coordinates {(1, -3658) (10, -5331) (100, -3676) (1000, -3368)};
\addplot[green!60!black, thick, mark=triangle*] coordinates {(1, -3082) (10, -1499) (100, -3980) (1000, -3605)};
\addplot[blue, thick, dashed, mark=o] coordinates {(1, 1528) (10, 297) (100, 0) (1000, 0)};
\addplot[red, thick, dashed, mark=square] coordinates {(1, 0) (10, 0) (100, 0) (1000, 0)};
\addplot[green!60!black, thick, dashed, mark=triangle] coordinates {(1, 0) (10, 0) (100, 0) (1000, 0)};
\legend{
    QALS ($k=1$),
    QALS ($k=10$),
    QALS ($k=100$),
    NO QALS ($k=1$),
    NO QALS ($k=10$),
    NO QALS ($k=100$)
}
\end{axis}
\end{tikzpicture}
\caption{Andamento di \textit{Profit GAP} al variare di TIMES per diversi valori di $k$.}
\label{fig:profit-gap}
\end{figure}

\begin{figure}[H]
\centering
\begin{tikzpicture}
\begin{axis}[
    width=0.95\textwidth,
    height=0.48\textwidth,
    xmode=log,
    log basis x=10,
    xlabel={TIMES},
    ylabel={Weight GAP},
    title={Weight GAP: Differenza tra la soluzione con fq minima trovata e la soluzione di DP},
    grid=both,
    legend style={at={(0.5,-0.25)}, anchor=north, legend columns=3},
    xtick={1,10,100,1000},
    xticklabels={1,10,100,1000},
    tick label style={font=\small},
    label style={font=\small},
    title style={font=\small},
]
\addplot[blue, thick, mark=*] coordinates {(1, -8233) (10, -6560) (100, -5652) (1000, -4434)};
\addplot[red, thick, mark=square*] coordinates {(1, -7658) (10, -6981) (100, -6131) (1000, -5382)};
\addplot[green!60!black, thick, mark=triangle*] coordinates {(1, -9475) (10, -6237) (100, -5371) (1000, -4255)};
\addplot[blue, thick, dashed, mark=o] coordinates {(1, -54) (10, -54) (100, 0) (1000, 0)};
\addplot[red, thick, dashed, mark=square] coordinates {(1, 0) (10, 0) (100, 0) (1000, 0)};
\addplot[green!60!black, thick, dashed, mark=triangle] coordinates {(1, 0) (10, 0) (100, 0) (1000, 0)};
\legend{
    QALS ($k=1$),
    QALS ($k=10$),
    QALS ($k=100$),
    NO QALS ($k=1$),
    NO QALS ($k=10$),
    NO QALS ($k=100$)
}
\end{axis}
\end{tikzpicture}
\caption{Andamento di \textit{Weight GAP} al variare di TIMES per diversi valori di $k$.}
\label{fig:weight-gap}
\end{figure}

\begin{figure}[H]
\centering
\begin{tikzpicture}
\begin{axis}[
    width=0.95\textwidth,
    height=0.48\textwidth,
    xmode=log,
    log basis x=10,
    xlabel={TIMES},
    ylabel={n of fQ Min Found},
    title={n of fQ Min Found},
    grid=both,
    legend style={at={(0.5,-0.25)}, anchor=north, legend columns=3},
    xtick={1,10,100,1000},
    xticklabels={1,10,100,1000},
    tick label style={font=\small},
    label style={font=\small},
    title style={font=\small},
]
\addplot[blue, thick, mark=*] coordinates {(1, 1) (10, 1) (100, 1) (1000, 1)};
\addplot[red, thick, mark=square*] coordinates {(1, 1) (10, 1) (100, 1) (1000, 1)};
\addplot[green!60!black, thick, mark=triangle*] coordinates {(1, 1) (10, 1) (100, 1) (1000, 1)};
\addplot[blue, thick, dashed, mark=o] coordinates {(1, 1) (10, 1) (100, 13) (1000, 94)};
\addplot[red, thick, dashed, mark=square] coordinates {(1, 1) (10, 8) (100, 65) (1000, 605)};
\addplot[green!60!black, thick, dashed, mark=triangle] coordinates {(1, 1) (10, 10) (100, 100) (1000, 1000)};
\legend{
    QALS ($k=1$),
    QALS ($k=10$),
    QALS ($k=100$),
    NO QALS ($k=1$),
    NO QALS ($k=10$),
    NO QALS ($k=100$)
}
\end{axis}
\end{tikzpicture}
\caption{Andamento di \textit{n of fQ Min Found} al variare di TIMES per diversi valori di $k$.}
\label{fig:n-of-fq-min-found}
\end{figure}

\end{document}